\section{Methodology}
\label{sec:method}
% [8 pts] Implementation details, key equations, architecture choices.

\subsection{Flow Matching}
\label{sec:fm}
Let $x_0\sim\mathcal N(0,I)$ be noise and $x_1$ a data sample.
FM defines a conditional path $\psi_t(x_0\mid x_1)$ between them and trains $v_\theta$ to match the path's velocity $u_t=\tfrac{d}{dt}\psi_t$ (conditional FM loss, Eq.~23 of~\cite{lipman2023flow}):
\begin{equation}
  \mathcal L_{\text{CFM}} = \mathbb E_{t,x_0,x_1}\big\|v_\theta(\psi_t,t) - u_t\big\|^2 .
  \label{eq:cfm}
\end{equation}
Although each conditional path is known, the loss has the same gradient as regressing onto the unknown marginal velocity field, so $v_\theta$ learns $\mathbb E[u_t\mid x_t]$.
I use straight paths:
\begin{itemize}
  \item \textbf{Task 1} (OT path): $\psi_t=(1-(1-\sigma_{\min})t)\,x_0+t\,x_1$, $u_t=x_1-(1-\sigma_{\min})x_0$, with $\sigma_{\min}=0.01$.
  \item \textbf{Tasks 2--4} (rectified path, as specified for the image tasks): $x_t=(1-t)x_0+t\,x_1$, $u_t=x_1-x_0$.
\end{itemize}
$t$ is drawn uniformly from $\{0,\tfrac{1}{1000},\dots,\tfrac{999}{1000}\}$ and $x_0,x_1$ are paired at random.
The image loss (\texttt{image\_common/fm.py}) is a direct translation of \cref{eq:cfm}:
\begin{lstlisting}
def get_loss(self, x1, class_label=None, x0=None):
    ...  # t ~ U, x0 ~ N(0, I) (provided)
    t_ = t.view(-1, *([1] * (x1.dim() - 1)))
    x_t = (1 - t_) * x0 + t_ * x1      # point on the straight path
    u_t = x1 - x0                      # constant target velocity
    model_out = self.network(x_t, t, class_label=class_label)
    return (model_out - u_t).pow(2).mean()
\end{lstlisting}
Task 1 is identical except for the OT path in \texttt{compute\_psi\_t} and the target $u_t$:
\begin{lstlisting}
psi_t = (1 - (1 - self.sigma_min) * t) * x + t * x1   # compute_psi_t
ut = x1 - (1 - self.fm_scheduler.sigma_min) * x0       # get_loss
\end{lstlisting}

\paragraph{Sampling.}
Starting from $x_0\sim\mathcal N(0,I)$, the ODE $\dot x=v_\theta(x,t)$ is integrated with $N$ uniform Euler steps, $x_{t+\Delta t}=x_t+\Delta t\,v_t$, ending at $t=1$.

\paragraph{Classifier-free guidance.}
For class-conditional generation the UNet receives a class label $c\in\{1,2,3\}$ (cat, dog, wild), whose embedding is added to the time embedding.
During training the label is replaced by a null label $0$ with probability 0.1, so a single network learns both $v(x,t\mid c)$ and $v(x,t\mid\varnothing)$.
At sampling time the two are combined as
\begin{equation}
  v_t = v(x_t,t\mid\varnothing) + w\,\big(v(x_t,t\mid c)-v(x_t,t\mid\varnothing)\big),
  \label{eq:cfg}
\end{equation}
where $w>1$ pushes samples further toward the class.
The sampling loop combines the Euler step and CFG:
\begin{lstlisting}
for i in range(num_inference_timesteps):  # t: 0 -> 1
    ...
    dt = expand_t(t_next - t, xt)
    if do_classifier_free_guidance:
        null_labels = torch.zeros_like(class_label)
        v_cond = self.network(xt, t, class_label=class_label)
        v_uncond = self.network(xt, t, class_label=null_labels)
        vt = v_uncond + guidance_scale * (v_cond - v_uncond)
    else:
        vt = self.network(xt, t, class_label=class_label)
    xt = self.fm_scheduler.step(xt, vt, dt)   # xt + dt * vt
\end{lstlisting}

\paragraph{Task 1 network.}
The velocity network reuses the Lab 1 \texttt{SimpleNet}: three time-conditioned linear layers (each output is scaled by an embedding of $t$) with ReLU, followed by a linear head.
Lab 1 fed integer timesteps to the sinusoidal embedding, whereas FM uses $t\in[0,1]$; I multiply $t$ by 1000 before embedding so the embedding varies over its intended range.
\begin{lstlisting}
def forward(self, x, t):
    t = t * self.num_timesteps   # [0, 1] -> [0, 1000)
    for layer in self.layers:    # TimeLinear x3
        x = F.relu(layer(x, t))
    return self.out_layer(x)
\end{lstlisting}

\subsection{Rectified Flow (Task 3)}
\label{sec:rf}
Because $x_0$ and $x_1$ are paired at random, straight conditional paths cross each other, and the averaged field $\mathbb E[u_t\mid x_t]$ produces \emph{curved} trajectories that need many Euler steps.
Reflow~\cite{liu2023rectified} replaces the random pairing with the deterministic pairing defined by the trained model:
\begin{enumerate}
  \item sample $x_0\sim\mathcal N(0,I)$ and a label $c$, and integrate the Task 2 model (20 steps, CFG 7.5) to obtain $z_1=\Phi(x_0)$;
  \item store $N=50{,}000$ triplets $(x_0,z_1,c)$;
  \item train a new model from scratch with the same loss (\cref{eq:cfm}) on $x_t=(1-t)x_0+t z_1$, $u_t=z_1-x_0$.
\end{enumerate}
ODE trajectories never cross, so the new pairs conflict far less and the retrained (1-RF) field is much straighter while generating the same distribution.
The only implementation detail that matters is to start the ODE from the \emph{saved} $x_0$ (passed explicitly to the sampler); otherwise the pairing is random again.
\begin{lstlisting}
# generate_reflow_data.py (x_0 and labels sampled above)
z_1 = fm.sample(
    shape,
    num_inference_timesteps=args.num_inference_steps,
    class_label=labels,
    guidance_scale=args.cfg_scale if args.use_cfg else 1.0,
    x0=x_0,   # keep the (x_0, z_1) coupling
)
\end{lstlisting}

\subsection{InstaFlow (Task 4)}
\label{sec:instaflow}
If trajectories were perfectly straight, one Euler step $x_1=x_0+v(x_0,0)$ would be exact.
InstaFlow trains a student for exactly this one-step map (Eq.~6 of~\cite{liu2024instaflow}):
\begin{equation}
  \mathcal L = \big\|x_0 + v_\theta(x_0,0\mid c) - x_1\big\|^2 + \mathcal L_{\text{LPIPS}},
\end{equation}
where $(x_0,x_1)$ are 50k pairs generated by the 1-RF teacher (20 steps, CFG $\alpha_2=1.5$) and LPIPS~\cite{zhang2018lpips} is an optional perceptual loss.
Guidance is baked into the targets, so sampling needs a single network evaluation without CFG.
\begin{lstlisting}
# get_loss
t_zero = torch.zeros(x0.shape[0], device=self.device)
v_pred = self.network(x0, t_zero, class_label=class_label)
x1_pred = x0 + v_pred                # one Euler step
loss = (x1_pred - x1).pow(2).mean()
if self.use_lpips:
    loss = loss + self.lpips_fn(x1_pred.clamp(-1, 1), x1.clamp(-1, 1)).mean()

# sample: one network call, no CFG
x_1 = x_T + self.network(x_T, t_zero, class_label=class_label)
\end{lstlisting}
Following the InstaFlow paper, I initialize the student from the 1-RF teacher instead of from scratch (\cref{sec:modified}).
The full pipeline is: FM (CFG 7.5) $\rightarrow$ 1-RF $\rightarrow$ InstaFlow (CFG 1.5) $\rightarrow$ one step.

\begin{table}[t]
  \centering
  \footnotesize
  \setlength{\tabcolsep}{3pt}
  \caption{\textbf{Training configuration.} Task 1 uses a 3-layer MLP of width 128; Tasks 2--4 share the same UNet (channels 128, multipliers $[1,2,2,2]$, attention at the 2nd level, 4 res-blocks per level). All runs use Adam on a single RTX 4090.}
  \label{tab:hparams}
  \begin{tabular}{@{}lccc@{}}
    \toprule
    & Task 1 & Task 2 / 3 & Task 4 \\
    \midrule
    Network     & MLP & UNet & UNet \\
    Init.       & random & random & 1-RF \\
    Batch       & 128 & 16 & 16 \\
    Steps       & 5k & 100k & 100k \\
    LR          & $5\!\times\!10^{-4}$ & $2\!\times\!10^{-4}$ & $2\!\times\!10^{-4}$ \\
    Warm-up     & -- & 200 & 5k + cosine \\
    CFG dropout & -- & 0.1 & 0 \\
    Loss        & CFM & CFM & L2 + LPIPS \\
    \bottomrule
  \end{tabular}
\end{table}


\subsection{Implementation Details}
\Cref{tab:hparams} lists the configurations, which follow the provided defaults.
Image quality is measured by FID~\cite{heusel2017fid} between 500 generated samples (random classes) and the 1,500 AFHQ validation images with the provided AFHQ Inception network.
I used PyTorch 2.8, so checkpoints, which store the whole network object, were loaded with \texttt{TORCH\_FORCE\_NO\_WEIGHTS\_ONLY\_LOAD=1}.
